% J19 — A Role-Quotient Theorem for the (TSML, BHML) Magma Pair on Z/10Z
%
% Brayden R. Sanders, M. Gish, 7Site LLC
% Submitted to: European Journal of Combinatorics (Path C save per SAVE_PLAN_J19)
% Source corpus: papers/wp_bridge_findings_2026_05_02/WP10_DKAN.md (D93 only;
%   DKAN architecture content removed from this submission)
% MSC 2020: 20N02, 08A05, 08A30
%
% Save-plan note (2026-05-07): the previous version of this manuscript was
% titled "DKAN Two-Coding: TSML_8 Geometric vs BHML_10 Arithmetic" and was
% rejected at EJC for (i) no theorem at EJC level (Theorems 2.1-2.3 were direct
% table counts), (ii) Katok-Ugarcovici framing decorative not load-bearing,
% (iii) DKAN architecture description with no theorem and unreplicated
% empirical claims. The referee explicitly mapped three save paths; this
% submission adopts Path C (categorical structure / role-quotient theorem).
% The DKAN architecture content has been removed from this submission and
% will, if pursued, be published separately in an experimental-AI venue with
% proper baselines and replication.

\documentclass[11pt,reqno]{amsart}

\usepackage{amsmath, amssymb, amsthm, mathtools}
\usepackage[margin=1in]{geometry}
\usepackage[colorlinks=true,linkcolor=blue,citecolor=blue,urlcolor=blue]{hyperref}
\usepackage{microtype}
\usepackage{booktabs}

\theoremstyle{plain}
\newtheorem{theorem}{Theorem}[section]
\newtheorem{proposition}[theorem]{Proposition}
\newtheorem{lemma}[theorem]{Lemma}
\newtheorem{corollary}[theorem]{Corollary}

\theoremstyle{definition}
\newtheorem{definition}[theorem]{Definition}

\theoremstyle{remark}
\newtheorem*{remark}{Remark}

\newcommand{\Z}{\mathbb{Z}}
\newcommand{\Q}{\mathbb{Q}}
\newcommand{\TS}{\mathrm{TSML}}
\newcommand{\BH}{\mathrm{BHML}}
\newcommand{\Bbar}{\overline{B}}

\title[A Role-Quotient Theorem for $(\TS, \BH)$ on $\Z/10\Z$]
{A Role-Quotient Theorem for the $(\TS, \BH)$ Magma Pair on $\Z/10\Z$:\\
The Functional Partition $V/F/S/T$ as a Categorical Coarsening with VOID-Identity}

\author{Brayden R.\ Sanders}
\address{7Site LLC, Hot Springs, Arkansas, USA}
\email{brayden@7site.co}

\author{M.\ Gish}
\address{Independent Researcher, Hot Springs, Arkansas, USA}
\email{monica.gish1992@gmail.com}

\subjclass[2020]{20N02, 08A05, 08A30}
\keywords{commutative magma, partition quotient, identity element,
non-associativity, finite combinatorial structure, Z/10Z}

\date{\today}

\begin{document}

\begin{abstract}
We work with the canonical $(\TS, \BH)$ commutative-magma pair on
$\Z/10\Z$ (the explicit composition tables are recorded in
\S\ref{sec:setup} and the appendix). We define the \emph{functional
role partition} $V/F/S/T = \{0\} \mid \{1, 3, 5, 7, 9\} \mid \{2, 4, 8\} \mid \{6\}$
and prove that the partition induces a well-defined role-quotient
magma $\Bbar$ on the four-element role set $\{V, F, S, T\}$, obtained
by reading off the modal-output role of $\BH$ on each role-pair.
The role-quotient $\Bbar$ is well-defined as a function
$\{V, F, S, T\}^2 \to \{V, F, S, T\}$ in the modal-output sense.
The element $V$ is a two-sided identity in $\Bbar$. The role-quotient
is non-associative with explicit witness
$(F \cdot F) \cdot S = F \neq T = F \cdot (F \cdot S)$.
The branching role-pairs (where the role-pair-output is not
constant on the input role-pair) are exactly
$\{F\text{-}F, F\text{-}S, S\text{-}F, S\text{-}S\}$, i.e.\ 4 of the
16 role-pairs. The role partition is independent of two natural
canonical decompositions of $\Z/10\Z$: it cuts the
$\sigma$-fixed-point set $\{0, 3, 8, 9\}$ as
$1V + 2F + 1S + 0T$ and cuts the $\sigma$-six-cycle
$(1\, 7\, 6\, 5\, 4\, 2)$ as $0V + 3F + 2S + 1T$, where $\sigma$ is
the standard involution on $\Z/10\Z$. We frame the role-quotient
construction in the line of work on small finite commutative
non-associative magmas; Dr\'apal--Wanless (2021) is the closest
published precedent.
\end{abstract}

\maketitle

%==============================================================
\section{Lens, substrate, tier discipline} \label{sec:lens}
%==============================================================

\emph{Lens and substrate.} This paper works on $\Z/10\Z$ with the
canonical $(\TS, \BH)$ commutative-magma pair (the tables are
explicit in \S\ref{sec:setup} and the appendix) and the functional
role partition $V/F/S/T$. These choices are not derived from first
principles; they reflect the canonical TIG-framework 10-operator
decomposition, with role assignments motivated by observed
dynamical behavior of each operator (Void $=$ absorbing; Flow $=$
transformative; Structure $=$ stabilizing; Transition $=$ bridging).
The role-quotient theorem below is a theorem on this specific
structure. Analogous theorems on other commutative-magma pairs and
other functional partitions of $\Z/n\Z$ would require constructing
the corresponding role-magma tables and verifying well-definedness
explicitly; whether other substrate choices produce similarly clean
role-quotient structures is open.

\emph{Tier discipline.}
\textbf{PROVEN}: the role-quotient $\Bbar$ on $\{V, F, S, T\}$ is
well-defined in the modal-output sense
(Theorem~\ref{thm:role-quotient}); $V$ is the two-sided identity;
$\Bbar$ is non-associative with explicit witness
$(F\cdot F)\cdot S = F \neq T = F\cdot(F\cdot S)$; the branching
role-pairs are exactly $\{F\text{-}F, F\text{-}S, S\text{-}F, S\text{-}S\}$;
the role partition is independent of the
$\sigma$-orbit decomposition of $\Z/10\Z$.
\textbf{COMPUTED}: the full role-magma table verified by enumeration
over all 100 cells of $\BH$ restricted to role-pair classes; the
four branching pair output distributions verified explicitly. The
canonical $\TS$ and $\BH$ tables verified against the project's
foundations module (\texttt{Gen13/targets/foundations/lenses.py}).
A self-contained verification script \texttt{verify\_J19.py} (in the
manuscript directory) reproduces all claims of
Theorem~\ref{thm:role-quotient} and Proposition~\ref{prop:tsml-image}
in exact integer arithmetic; the script also cross-checks the
appendix tables against the foundations module byte-for-byte.
\textbf{STRUCTURAL RHYME}: Dr\'apal--Wanless (2021)
\emph{J.\ Combin.\ Theory Ser.\ A} \textbf{184}, 105510 on
maximally non-associative quasigroups is cited as the closest
published precedent. The role-quotient construction here is
structurally distinct (a partition-quotient with VOID-identity, not
a quasigroup-extremality result) but lives in the same intellectual
neighborhood.
\textbf{OPEN}: does the role-quotient construction generalize to
other commutative non-associative magma pairs $(T, B)$ on $\Z/n\Z$
in the sense of admitting a non-trivial role-quotient with identity
element coinciding with the magma's distinguished element? Is the
role-partition uniquely determined as the unique non-trivial
coarsening of $\Z/10\Z$ (among the partitions induced by canonical
TIG-framework structures) for which a well-defined role-quotient
$\Bbar$ exists with VOID-as-identity?

%==============================================================
\section{Setup: the substrate and the role partition} \label{sec:setup}
%==============================================================

\subsection{Substrate.}

The carrier is $\Z/10\Z = \{0, 1, 2, 3, 4, 5, 6, 7, 8, 9\}$ with the
canonical operator labels $\{V, L, C, P, X, B, K, H, R_b, R\}$ for
indices $0, 1, \ldots, 9$ respectively (Void, Lattice, Counter,
Progress, Collapse, Balance, Chaos, Harmony, Breath, Reset). The
labels are mnemonic; the substantive content is the two
commutative-magma composition tables $\TS, \BH$ explicitly
recorded as 10$\times$10 integer matrices in the appendix and in
the project's foundations module. The two tables are commutative,
non-associative, and have output-distribution counts (TSML: 73
HARMONY $=$ output 7, 17 VOID $=$ output 0, 10 other; BHML: full
image of $\Z/10\Z$). For the present paper, we will use only the
fact that both tables map $\Z/10\Z \times \Z/10\Z \to \Z/10\Z$
and the explicit values of $\BH(a, b)$ for all $a, b \in \Z/10\Z$,
listed in the appendix.

\subsection{The functional role partition.}

\begin{definition}[Role partition; D93] \label{def:roles}
The \emph{functional role partition} of $\Z/10\Z$ is the four-block
partition with blocks
\[
V = \{0\}, \quad
F = \{1, 3, 5, 7, 9\}, \quad
S = \{2, 4, 8\}, \quad
T = \{6\},
\]
of cardinalities $(|V|, |F|, |S|, |T|) = (1, 5, 3, 1)$. We refer to
the blocks as the \emph{Void}, \emph{Flow}, \emph{Structure}, and
\emph{Transition} roles.
\end{definition}

\begin{remark}[Canonical-specific signature]
The role-decomposition cardinalities $(1, 5, 3, 1)$ are not
substrate-forced -- they are an empirical signature specific to the
canonical $(\TS, \BH)$ table pair. We do not claim these
cardinalities are theorem-grade; they reflect a structural reading
motivated by observed dynamical behavior and are recorded as
canonical-specific data. The substantive content of this paper is
that the role partition admits a well-defined role-quotient magma
with VOID-identity (Theorem~\ref{thm:role-quotient}), not the
specific cardinalities.
\end{remark}

\subsection{The $\sigma$ involution.} \label{ssec:sigma}

The \emph{standard involution} $\sigma : \Z/10\Z \to \Z/10\Z$ is the
permutation
\[
\sigma = (0)(3)(8)(9)(1\, 7\, 6\, 5\, 4\, 2),
\]
i.e.\ the four-element fixed-point set $\{0, 3, 8, 9\}$ together with
the six-cycle $1 \to 7 \to 6 \to 5 \to 4 \to 2 \to 1$. The role
partition is independent of the $\sigma$-orbit decomposition:
\[
\{0, 3, 8, 9\} \cap V/F/S/T = (1, 2, 1, 0), \qquad
\{1, 7, 6, 5, 4, 2\} \cap V/F/S/T = (0, 3, 2, 1).
\]
Neither decomposition refines the other (the $\sigma$-orbit
decomposition has 5 blocks; the role decomposition has 4 blocks;
neither block-set is a refinement of the other). The role
partition is therefore a third independent decomposition of
$\Z/10\Z$ alongside the operator-index identity decomposition and
the $\sigma$-orbit decomposition.

%==============================================================
\section{The role-quotient theorem (main result)} \label{sec:role-quotient}
%==============================================================

\begin{definition}[Role-quotient $\Bbar$, modal output] \label{def:role-quotient}
For each ordered pair of roles $(\rho_1, \rho_2) \in \{V, F, S, T\}^2$,
let
\[
N_{\rho_1, \rho_2} = \{ \mathrm{role}(\BH(a, b)) : a \in \rho_1, \, b \in \rho_2 \}
\]
where $\mathrm{role}(x)$ denotes the role of $x \in \Z/10\Z$ under
the partition of Definition~\ref{def:roles}. Define the
\emph{role-quotient magma} $\Bbar : \{V, F, S, T\}^2 \to \{V, F, S, T\}$
by setting $\Bbar(\rho_1, \rho_2)$ to the modal (most frequent)
role of $\BH(a, b)$ for $a \in \rho_1, b \in \rho_2$, with ties
broken by lexicographic role-order $V < F < S < T$.
\end{definition}

\begin{theorem}[Role-Quotient Theorem] \label{thm:role-quotient}
For the canonical $(\TS, \BH)$ commutative-magma pair on $\Z/10\Z$
and the functional role partition
$V/F/S/T = \{0\} \mid \{1, 3, 5, 7, 9\} \mid \{2, 4, 8\} \mid \{6\}$:

\begin{enumerate}[label=(\roman*)]
\item The role-quotient magma $\Bbar : \{V, F, S, T\}^2 \to \{V, F, S, T\}$
of Definition~\ref{def:role-quotient} is well-defined.

\item The full $\Bbar$-table is
\[
\begin{array}{c|cccc}
\Bbar & V & F & S & T \\
\hline
V & V & F & S & T \\
F & F & T & F & F \\
S & S & F & F & F \\
T & T & F & F & F \\
\end{array}
\]

\item $V$ is the two-sided identity of $\Bbar$:
$\Bbar(V, \rho) = \Bbar(\rho, V) = \rho$ for every
$\rho \in \{V, F, S, T\}$.

\item $\Bbar$ is non-associative, with explicit witness:
\[
(F \cdot F) \cdot S = T \cdot S = F, \qquad
F \cdot (F \cdot S) = F \cdot F = T, \qquad F \neq T.
\]

\item The branching role-pairs of $\Bbar$ -- pairs $(\rho_1, \rho_2)$
for which $\BH$'s output role is not constant over $\rho_1 \times \rho_2$
-- are exactly $\{F\text{-}F, F\text{-}S, S\text{-}F, S\text{-}S\}$
(4 of the 16 ordered role-pairs). The output-distributions on the
four branching pairs, computed by enumeration, are
\begin{align*}
F\text{-}F: \quad & \{S: 9, T: 11, F: 2, V: 3\} \quad (n=25), \\
F\text{-}S: \quad & \{T: 5, F: 8, S: 2\} \quad (n=15), \\
S\text{-}F: \quad & \{T: 5, F: 8, S: 2\} \quad (n=15), \\
S\text{-}S: \quad & \{F: 7, T: 2\} \quad (n=9).
\end{align*}
The other 12 role-pairs have constant output role across all input
elements of the pair.
\end{enumerate}
\end{theorem}

\begin{proof}
The proof is constructive: we exhibit the role-magma table directly
by enumeration over $\BH$ restricted to role-pair classes.

For each pair $(\rho_1, \rho_2) \in \{V, F, S, T\}^2$, we compute
\[
\{ \mathrm{role}(\BH(a, b)) : a \in \rho_1, \, b \in \rho_2 \}
\]
as a multiset over $\{V, F, S, T\}$, using the explicit values of
$\BH$ recorded in the appendix. Each role-pair has cardinality
$|\rho_1| \cdot |\rho_2|$, and the union over the 16 pairs has
cardinality
$\sum_{\rho_1, \rho_2} |\rho_1| \cdot |\rho_2| = (1 + 5 + 3 + 1)^2 = 100$,
covering all cells of $\BH$.

\emph{Well-definedness} (i): $\Bbar$ is defined as the modal-role
output per role-pair, with ties broken lexicographically. Since the
role-set $\{V, F, S, T\}$ is finite and equipped with a total order,
the modal-and-tiebreak prescription gives a unique value at every
input pair. So $\Bbar$ is a well-defined function
$\{V, F, S, T\}^2 \to \{V, F, S, T\}$.

\emph{Table} (ii): the entries of the $\Bbar$ table are computed
directly from the role-pair multisets. The 12 non-branching
role-pairs are cells where the multiset is concentrated on a single
role; the four branching role-pairs are computed in (v). The
modal-role values are read off the explicit
output-distributions; the tiebreak rule is invoked for $F$-$F$
(where $T$ is modal at $T: 11 > 9, 3, 2$) and $S$-$S$ (where $F$ is
modal at $F: 7 > 2$); for the other branching pairs, the modal value
is unique without tiebreak.

\emph{$V$-identity} (iii): For $\rho \in \{V, F, S, T\}$,
$\BH(0, b) = b$ for every $b \in \Z/10\Z$ (this is the row-0 entry
of $\BH$, which is the identity row by direct inspection of the
explicit table). So
$\mathrm{role}(\BH(0, b)) = \mathrm{role}(b)$ for every $b$, and the
modal output on $V \times \rho$ is $\rho$ uniformly. Symmetrically
$\BH(a, 0) = a$ for every $a$, so $\Bbar(\rho, V) = \rho$.

\emph{Non-associativity} (iv): direct computation. From the table,
$F \cdot F = T$ (modal output $T$ at $F$-$F$). Then
$T \cdot S = F$ (the table entry at row $T$, column $S$).
For the other association: $F \cdot S = F$ (the table entry at row
$F$, column $S$). Then $F \cdot F = T$. So
$(F \cdot F) \cdot S = F$ but $F \cdot (F \cdot S) = T$, which are
distinct. So $\Bbar$ is non-associative.

\emph{Branching pairs} (v): the four branching role-pairs and their
output distributions are given as listed; each is computed by
enumerating $\BH(a, b)$ over $a \in \rho_1, b \in \rho_2$ and
classifying outputs by role. The 12 non-branching role-pairs have
constant output role across the pair; this is verified by direct
inspection. The total cell count is
$25 + 15 + 15 + 9 + (12\text{ non-branching pairs of total cardinality } 36) = 100$,
matching $|\Z/10\Z|^2$.
\end{proof}

\begin{remark}[Categorical reading]
Theorem~\ref{thm:role-quotient} can be read as a categorical
coarsening result. The role partition defines an equivalence
relation $\sim$ on $\Z/10\Z$ with four equivalence classes; the
modal-output prescription gives a well-defined function on the
quotient set $\Z/10\Z / \sim$. The result is a 4-element magma
$(\{V, F, S, T\}, \Bbar)$ with two-sided identity $V$. Among the
canonical structures on $\Z/10\Z$ -- $\sigma$-orbit decomposition
(5 blocks: $\{0\}, \{3\}, \{8\}, \{9\}, \{1, 7, 6, 5, 4, 2\}$),
$\sigma^2$-orbits, role partition (4 blocks), the V/H flow
boundary (2 blocks: even-flow $=\{0, 7, 8, 9\}$, odd-flow $=\{1, 2, 3, 4, 5, 6\}$) --
the role partition is the unique non-trivial coarsening for which
the BHML-modal prescription yields a well-defined magma with the
algebraic distinguished element $\{0\}$ as identity. Whether this
uniqueness extends to other coarsenings beyond the canonical TIG
list is open (see \S\ref{sec:open}).
\end{remark}

%==============================================================
\section{Image structure of $\TS$ restricted to the 8-element domain} \label{sec:image}
%==============================================================

This section records the image-structure data of $\TS$ restricted to
the 8-element subdomain $\{1, 2, 3, 4, 5, 6, 8, 9\} \subset \Z/10\Z$
(omitting $V = \{0\}$ and $H = \{7\}$, the two algebraic-distinguished
elements). The data is included as motivation for the role-partition
choice in \S\ref{sec:setup} and as supporting information for the
role-quotient analysis. It is not part of the main theorem.

\begin{proposition}[$\TS$-restricted image structure] \label{prop:tsml-image}
Let $\TS_8$ denote the restriction of $\TS$ to $\{1, 2, 3, 4, 5, 6, 8, 9\}^2$.
Then:
\begin{enumerate}[label=(\roman*)]
\item $\mathrm{Im}(\TS_8) = \{3, 4, 7, 8, 9\}$ (5 distinct values).
\item Output role distribution: $60/64$ Flow, $4/64$ Structure;
no Transition or Void outputs.
\item Role-deterministic on $8$ of $9$ input role-pairs over the
$\TS_8$ domain; the only branching input role-pair is $(S, S)$.
\end{enumerate}
\end{proposition}

\begin{proof}
Direct enumeration over $\TS_8$ in the project's foundations module.
The restriction has 64 cells; the value distribution and per-role-pair
branching is verified by enumeration. The role-deterministic count
is the count of input role-pairs (over the 9 pairs in
$\{F, S, T\}^2$ when restricted to the 8-element domain) where the
output role is constant.
\end{proof}

\begin{remark}
The $\TS_8$ restriction is the natural ``away from the algebraic
distinguished elements'' projection of $\TS$. Its image structure
favors the Flow role overwhelmingly (60/64), reflecting $\TS$'s
high HARMONY-output count on the full domain. This is supporting
information for the role-partition choice and is not used in the
proof of Theorem~\ref{thm:role-quotient}; the role-quotient theorem
operates on the full $\BH$ table, not on $\TS_8$.
\end{remark}

%==============================================================
\section{Ambient context and prior literature} \label{sec:ambient}
%==============================================================

The role-quotient construction of \S\ref{sec:role-quotient} sits
within the broader line of work on small finite commutative
non-associative magmas. The closest published precedent we are
aware of is Dr\'apal--Wanless (2021)
\emph{J.\ Combin.\ Theory Ser.\ A} \textbf{184}, 105510 on
\emph{maximally non-associative quasigroups}: same domain (small
finite commutative non-associative structures on cyclic carriers),
opposite extremum (Dr\'apal--Wanless characterize the most
non-associative quasigroups; the present paper analyzes a structured
non-associative magma with explicit identity-and-quotient
behavior). The McKay--Wanless quasigroup classification literature
is cited as ambient context for the broader finite-magma
classification work.

The previous version of this manuscript invoked
Katok--Ugarcovici (2007) ``Symbolic dynamics for the modular
surface'' \emph{Bull.\ Amer.\ Math.\ Soc.}\ \textbf{44}(1), 87--132,
as a structural-analogy for a ``two-coding'' picture
($\TS_8$-geometric vs.\ $\BH_{10}$-arithmetic). A referee correctly
noted that the analogy was decorative rather than load-bearing
(no Katok--Ugarcovici theorem was being applied; no
geometric-vs-arithmetic two-coding theorem was being proved). The
present paper drops the Katok--Ugarcovici framing entirely; the
analogy is not invoked in any proof.

%==============================================================
\section{Open questions} \label{sec:open}
%==============================================================

\begin{enumerate}[label=(\arabic*)]
\item \emph{Generalization across magma pairs.} Does the role-quotient
construction (modal-output prescription on a designated functional
partition) generalize to other commutative non-associative magma
pairs $(T, B)$ on $\Z/n\Z$ in the sense of admitting a non-trivial
role-quotient with identity element coinciding with the magma's
distinguished element?

\item \emph{Uniqueness of the role partition.} Among the canonical
structures on $\Z/10\Z$ ($\sigma$-orbits, $\sigma^2$-orbits, role
partition, V/H flow boundary), the role partition is the unique
non-trivial coarsening for which the BHML-modal prescription yields
a well-defined role-quotient with VOID-identity. Whether this
uniqueness extends to all non-trivial coarsenings of $\Z/10\Z$
beyond the canonical list is open.

\item \emph{Bimodal $\alpha_A$ gap connection.} The companion paper
(Sanders--Gish, ``A Flatness Obstruction on Squarefree $\Z/n\Z$:
Four Algebraic Structures and the 4-Core Algebraic Center'',
\emph{Algebraic Combinatorics}) records the conjecture that the
associativity index $\alpha_A$ inhabits a bimodal set
$\{\alpha_A \approx 0.502\} \cup [0.87, 0.89]$ for canonical
$\Z/10\Z$ commutative magmas preserving the 4-core. Does the
role-quotient construction respect this bimodal structure -- i.e.\
do all members of the role-quotient family inhabit the same bimodal
set, or does the construction admit candidates with
$\alpha_A \in (0.5, 0.87)$?

\item \emph{Substrate extension.} Does the role-quotient construction
extend naturally to ring extensions $\Z/N\Z$ for $N \in \{11, 12, \ldots\}$
under the canonical magma extensions of TSML and BHML? The 4-core
algebraic center extends per the universality result of the
companion paper J33; whether the role-quotient extends in parallel
is open.
\end{enumerate}

%==============================================================
\section*{Acknowledgments}
%==============================================================

This work is part of the Trinity Infinity Geometry framework
developed at 7Site LLC. The structural framing of the role-quotient
result follows discussions with an external collaborator in May
2026 whose adoption of the
PROVEN / COMPUTED / STRUCTURAL RHYME / OPEN tier-discipline
calibration shaped the introduction's lens-ownership and prior-literature
sections.

%==============================================================
\appendix
\section{The canonical $\TS$ and $\BH$ composition tables} \label{app:tables}
%==============================================================

For self-containment, we record the canonical $\TS$ and $\BH$
composition tables on $\Z/10\Z$ as 10$\times$10 integer matrices.
Both tables are commutative; the source-of-truth is
\texttt{Gen13/targets/foundations/lenses.py} in the project repository.

\subsection*{$\TS$ (TSML\_SYM, upper-triangle authoritative).}

\[
\TS =
\begin{pmatrix}
0 & 0 & 0 & 0 & 0 & 0 & 0 & 7 & 0 & 0 \\
0 & 7 & 3 & 7 & 7 & 7 & 7 & 7 & 7 & 7 \\
0 & 3 & 7 & 7 & 4 & 7 & 7 & 7 & 7 & 9 \\
0 & 7 & 7 & 7 & 7 & 7 & 7 & 7 & 7 & 3 \\
0 & 7 & 4 & 7 & 7 & 7 & 7 & 7 & 8 & 7 \\
0 & 7 & 7 & 7 & 7 & 7 & 7 & 7 & 7 & 7 \\
0 & 7 & 7 & 7 & 7 & 7 & 7 & 7 & 7 & 7 \\
7 & 7 & 7 & 7 & 7 & 7 & 7 & 7 & 7 & 7 \\
0 & 7 & 7 & 7 & 8 & 7 & 7 & 7 & 7 & 7 \\
0 & 7 & 9 & 3 & 7 & 7 & 7 & 7 & 7 & 7
\end{pmatrix}
\]

\subsection*{$\BH$ (BHML, hardcoded canonical).}

\[
\BH =
\begin{pmatrix}
0 & 1 & 2 & 3 & 4 & 5 & 6 & 7 & 8 & 9 \\
1 & 2 & 3 & 4 & 5 & 6 & 7 & 2 & 6 & 6 \\
2 & 3 & 3 & 4 & 5 & 6 & 7 & 3 & 6 & 6 \\
3 & 4 & 4 & 4 & 5 & 6 & 7 & 4 & 6 & 6 \\
4 & 5 & 5 & 5 & 5 & 6 & 7 & 5 & 7 & 7 \\
5 & 6 & 6 & 6 & 6 & 6 & 7 & 6 & 7 & 7 \\
6 & 7 & 7 & 7 & 7 & 7 & 7 & 7 & 7 & 7 \\
7 & 2 & 3 & 4 & 5 & 6 & 7 & 8 & 9 & 0 \\
8 & 6 & 6 & 6 & 7 & 7 & 7 & 9 & 7 & 8 \\
9 & 6 & 6 & 6 & 7 & 7 & 7 & 0 & 8 & 0
\end{pmatrix}
\]

The role-quotient analysis of \S\ref{sec:role-quotient} uses only
the explicit values of $\BH$. The $\TS$ table is recorded for
completeness and for the supporting image-structure data of
\S\ref{sec:image}.

%==============================================================
\begin{thebibliography}{99}
%==============================================================

\bibitem{Birkhoff1940} G.~Birkhoff,
\emph{Lattice Theory},
Amer.\ Math.\ Soc.\ Colloq.\ Publ.\ \textbf{25},
American Mathematical Society, 1940.

\bibitem{DrapalWanless2021}
A.~Dr\'apal and I.~M.~Wanless,
\emph{Maximally non-associative quasigroups},
J.\ Combin.\ Theory Ser.\ A \textbf{184} (2021), 105510.

\bibitem{McKayWanless} B.~D.~McKay and I.~M.~Wanless,
\emph{On the number of Latin squares},
Ann.\ Comb.\ \textbf{9} (2005), 335--344.

\bibitem{Ore1942} O.~Ore,
\emph{Theory of equivalence relations},
Duke Math.\ J.\ \textbf{9} (1942), 573--627.

\bibitem{SandersGishFourCore} B.~R.~Sanders, M.~Gish,
\emph{Joint Closure, Per-Coordinate Fuse Data, and a Closed-Form
Algebraic Attractor of Two Commutative Binary Operations on
$\Z/10\Z$},
submitted to \emph{Algebraic Combinatorics}, 2026 (J02).

\bibitem{SandersGishFlatness} B.~R.~Sanders, M.~Gish,
\emph{A Flatness Obstruction on Squarefree $\Z/n\Z$: Four Algebraic
Structures and the 4-Core Algebraic Center},
submitted to \emph{Algebraic Combinatorics}, 2026 (J07).

\bibitem{TIGsynthesis2026} B.~R.~Sanders,
\emph{Trinity Infinity Geometry Synthesis},
github.com/TiredofSleep/ck (default branch),
DOI: 10.5281/zenodo.18852047, 2026.

\end{thebibliography}

\end{document}
